\section{Universal reductions to binary-meet deletion}
\label{sec:universal-reductions}

This section distinguishes three logically equivalent \emph{universal theorem schemes}.  The equivalence does not mean that every individual nonmaximal deletion is itself a cover or a binary meet.

\subsection{Immediate covers}

\begin{lemma}[Finite lower intervals]
\label{lem:finite-lower-interval}
For every finite code $\C$, the interval
\[
 [\{\varnothing\},\C]=\{\D:\D\le\C\}/\cong
\]
is finite.
\end{lemma}

\begin{proof}
A nontrivial output coordinate of a surjective morphism from $\C$ is determined by the inverse image of its simple trunk, which is a proper trunk of $\C$.  After deleting trivial coordinates and identifying duplicate coordinates, an image is determined, up to permutation, by a finite set of distinct proper trunks.  If $t(\C)$ is the number of proper trunks, there are at most $2^{t(\C)}$ such sets.  Further removal of redundant coordinates can only decrease the count.
\end{proof}

\begin{lemma}[Morphic preservation]
\label{lem:morphic-preservation}
A surjective code morphic image of an open convex code is open convex.  A surjective morphic image of a polytope-convex code is polytope convex.
\end{lemma}

\begin{proof}
For each nontrivial output coordinate $j$, write the inverse image of its simple trunk as $\Tk_\C(\rho_j)$ and define
\[
 V_j=\bigcap_{i\in\rho_j}U_i.
\]
Then $x\in V_j$ exactly when the old word at $x$ belongs to the defining trunk, so the new word is the morphic image of the old word.  Intersections of open convex sets are open convex.  If the old fields are interiors of polytopes, replace $V_j$ by the interior of $\bigcap_{i\in\rho_j}P_i$; nonemptiness of the trunk ensures the usual interior-intersection identity.  The mandatory empty word belongs to no proper trunk.
\end{proof}

\begin{theorem}[Saturated-chain reduction]
\label{thm:saturated-chain}
Universal polyhedralization is equivalent to the following local theorem scheme:
\[
 \A\lessdot\B,
 \quad \A\text{ polytope convex},
 \quad \B\text{ open convex}
 \quad\Longrightarrow\quad
 \B\text{ polytope convex}.
 \tag{CL}
\]
\end{theorem}

\begin{proof}
Universal polyhedralization plainly implies (CL).  Conversely, the finite interval below an open convex code $\C$ contains a saturated chain
\[
 \{\varnothing\}=\C_0\lessdot\C_1\lessdot\cdots\lessdot\C_m=\C.
\]
Every $\C_s$ is a minor of $\C$ and hence open convex by \Cref{lem:morphic-preservation}.  The one-word code is polytopal.  Repeated application of (CL) proves every $\C_s$ polytopal.
\end{proof}

The covering relations in the code poset have now been classified by \citet{JeffsTrang2025}.  Our later reductions do not require the four-type classification; they pass instead through monotonicity and maximal-word intersections.

\subsection{Single nonmaximal deletion}

Let $\widehat\C$ denote the intersection completion generated by the maximal words of $\C$ together with all requested words; equivalently, it is the smallest intersection-complete code containing $\C$ without changing the maximal words.

\begin{theorem}[Single-codeword-deletion reduction]
\label{thm:single-deletion}
Universal polyhedralization is equivalent to the theorem scheme:
\begin{quote}
If $\D$ is polytope convex, $\E=\D\setminus\{\sigma\}$ is open convex, and $\sigma$ is nonmaximal in $\D$, then $\E$ is polytope convex. \hfill(DEL)
\end{quote}
\end{theorem}

\begin{proof}
Universal polyhedralization implies (DEL).  Conversely, fix an open convex $\C$.  By \Cref{thm:mic-polytopal}, its intersection completion $\widehat\C$ is polytope convex.  Delete the finitely many words of $\widehat\C\setminus\C$ one at a time.  Every intermediate code lies between $\C$ and $\widehat\C$ and has the same maximal words, so it is open convex by \Cref{cor:open-monotonicity}.  Every deleted word is nonmaximal because all maximal words already lie in $\C$.  Applying (DEL) at each step ends at a polytopal realization of $\C$.
\end{proof}

\subsection{Binary meets}

Define
\[
 \Bcore(\C)=\{\varnothing\}\cup
 \left\{\bigcap_{a\in R}m_a:
 \varnothing\ne R\subseteq[k]\right\},
\]
where $m_1,\ldots,m_k$ are the maximal words of $\C$, and put
\[
 \C^+=\C\cup\Bcore(\C).
\]

\begin{theorem}[Binary-meet deletion reduction]
\label{thm:binary-meet-reduction}
Let $\C$ be open convex.  Then $\C^+$ is polytope convex.  Moreover, the words of $\Bcore(\C)\setminus\C$ may be ordered as $\sigma_1,\ldots,\sigma_r$ so that, with
\[
 \D_0=\C^+,
 \qquad
 \D_s=\D_{s-1}\setminus\{\sigma_s\},
\]
every $\D_s$ is open convex and, at the $s$th deletion, there are strict surviving superwords $\tau_s,\upsilon_s\in\D_s$ with
\[
 \sigma_s=\tau_s\cap\upsilon_s.
\]
Consequently, universal polyhedralization is equivalent to proving all open-convex binary-meet deletions from polytopal lower codes are polytopal.
\end{theorem}

\begin{proof}
The core $\Bcore(\C)$ is polytope convex, and
\[
 \Bcore(\C)\subseteq\C^+\subseteq\Delta(\Bcore(\C)).
\]
Thus $\C^+$ is polytope convex by monotonicity.  It is also open convex because
\[
 \C\subseteq\C^+\subseteq\Delta(\C).
\]

Order the missing core words by nondecreasing cardinality, breaking ties arbitrarily.  Fix the current word $\sigma$.  Choose an inclusion-minimal nonempty family $R\subseteq\Max(\C)$ satisfying
\[
 \sigma=\bigcap_{m\in R}m.
\]
The family has at least two members, because a one-member intersection would be a maximal word already in $\C$.  Inclusion minimality implies
\[
 \bigcap_{m\in R'}m\supsetneq\sigma
\]
for every proper nonempty $R'\subsetneq R$.

Partition $R=A\sqcup B$ into two nonempty subfamilies and put
\[
 \tau=\bigcap_{m\in A}m,
 \qquad
 \upsilon=\bigcap_{m\in B}m.
\]
Then $\tau,\upsilon\supsetneq\sigma$ and $\tau\cap\upsilon=\sigma$.  Their cardinalities are strictly larger than $|\sigma|$.  If one is absent from the original code, it appears later in the nondecreasing-cardinality order and has not yet been deleted; if it lies in $\C$, it is never deleted.  Hence both survive the current deletion.

Every intermediate code lies between $\C$ and $\C^+$ with unchanged maximal words, so open-convex monotonicity applies.  A universal binary-meet deletion theorem therefore yields polytopality at every step.  The converse is immediate from universal polyhedralization.
\end{proof}

\begin{figure}[t]
\centering
\begin{tikzpicture}[>=Stealth,every node/.style={font=\small},word/.style={draw,ellipse,minimum width=18mm,minimum height=8mm}]
\node[word] (tau) {$\tau$};
\node[word,right=28mm of tau] (ups) {$\upsilon$};
\node[word,below=16mm of $(tau)!0.5!(ups)$] (sig) {$\sigma=\tau\cap\upsilon$};
\draw[->] (tau)--(sig);
\draw[->] (ups)--(sig);
\node[below=4mm of sig,font=\small,align=center] {delete $\sigma$; keep the two strict superwords};
\end{tikzpicture}
\caption{A binary-meet deletion.  The absent word $\sigma$ is the intersection of two strict surviving superwords $\tau$ and $\upsilon$.}
\label{fig:binary-meet}
\end{figure}

\begin{remark}[Logical strength]
The theorem is a reduction of universal statements.  It does not assert that an arbitrary individual nonmaximal word is a binary meet in the code from which it is deleted.  The binary-meet property is guaranteed only along the canonical cardinality-ordered deletion path from $\C^+$ to $\C$.
\end{remark}
